
% \begin{restatable*}{proposition}{optTruthfulVsOptUntruthful}
%     For any $k\geq 0$ there exists an environment such that  
%     \(
%         \EE\sbr{
%             \rewardss{
%                 \optstrategytt{
%                     \optstoptime
%                 }
%             }
%         }
%         -
%         \EE\sbr{
%             \rewardss{
%                 \opttruthstrategytt{
%                     \opttruthstoptime
%                 }
%             }
%         }
%         \geq k
%     \),
%     where 
%     $\optexpolicy$
%     and 
%     $\opttruthexpolicy$
%     are the optimal policies assuming and not assuming truthful reporting respectively.
% \end{restatable*}  


% \begin{restatable*}{proposition}{opt-truthful-vs-opt-untruthful2}
% % \label{prop:untruthfulreporting}
%     There exist a sequence of model instances such that 
%     \mycomment{typo, remove $k$, just one problem instance}
%     \[
%         \EE\sbr{
%             \rewardss{
%                 \optstrategytt{
%                     \optstoptime
%                 }
%             }
%         }
%         -
%         \EE\sbr{
%             \rewardss{
%                 \opttruthstrategytt{
%                     \opttruthstoptime
%                 }
%             }
%         }
%         \xrightarrow{k\rightarrow\infty}\infty
%     \]
%     where 
%     \(
%         \cbr{
%             \optexpolicy_k
%         }_{k\geq 0}
%         =
%         \cbr{
%             \rbr{
%                 \cbr{
%                     \optstrategyt
%                 }_{\timeind\geq 0},
%                 \optstoptime
%             }_k
%         }_{k\geq 0}
%     \)
%     and 
%     \(
%         \cbr{
%             \opttruthexpolicy_k
%         }_{k\geq 0}
%         =
%         \cbr{
%             \rbr{
%                 \cbr{
%                     \opttruthstrategyt
%                 }_{\timeind\geq 0},
%                 \opttruthstoptime
%             }_k
%         }_{k\geq 0}
%     \)
%     are sequences of optimal and optimal truthful policies for each model.
% \end{restatable*}  


% \begin{proof}
%     Let $\numdataa=\numdatab=1$.
%     For a particular $k\in\NN$, consider the valuations
%     % utilties, and welfare function,
%     % \(
%     % \)
%     \begin{align*}
%         \valaii{0}=\valbjj{0}=0
%         \qquad
%         \valaii{1}
%         =
%         \begin{cases}
%             2k
%             &
%             \text{w.p. } \frac{1}{2}
%             \\
%             % \frac{k}{2}
%             k
%             &
%             \text{w.p. } \frac{1}{2}
%         \end{cases}  
%         \qquad
%         \valbjj{1} = 2k
%     \end{align*}  
%     and let the competitors' utilities and the welfare function be given by
%     \begin{align*}
%         \utila(x,y) 
%         =
%         x-y
%         \qquad
%         \utilb(x,y)
%         =
%         \begin{cases}
%             0 
%             &
%             \text{if }
%             (y-x)\geq -k
%             %  k\geq x-y
%             \\
%             2(y-x)
%             &
%             \text{otherwise}
%         \end{cases}  
%         \qquad
%         \welfarefunc(x,y)=x+y .
%     \end{align*}
%     First consider when competitors are truthful. 
%     Under this model, the optimal policy 
%     \(
%         \optexpolicy_k
%         % =
%         % \rbr{
%         %     \cbr{
%         %         \optstrategyt
%         %     }_{\timeind\geq 0},
%         %     \optstoptime
%         % }_k
%     \)
%     first gives $\competitora$ competitor $\competitorb$'s point,
%     if 
%     % $\valaii{1}=\frac{k}{2}$ 
%     $\valaii{1}=k$ 
%     then sharing terminates, otherwise $\competitorb$ is given
%     $\competitora$'s point to ensure $\competitorb$'s utility does not lead to a negative welfare.
%     Thus the expected welfare under this policy is 
%     \[
%         \EE\sbr{
%             \rewardss{
%                 % \optstrategytt{
%                 %     \optstoptime_k
%                 % }^k
%                 \stopstateforpolicy{
%                     \optexpolicy_k
%                 }
%             }
%         }
%         =
%         \frac{1}{2}
%         \welfarefunc\rbr{
%             \utilass{
%                 (1,0)
%             },
%             \utilbss{
%                 (1,0)
%             }
%         }
%         +
%         \frac{1}{2}
%         \welfarefunc\rbr{
%             \utilass{
%                 (1,1)
%             },
%             \utilbss{
%                 (1,1)
%             }
%         }
%         =
%         \frac{1}{2}
%         \rbr{
%             % \frac{k}{2}
%             k
%             + 0 
%         }
%         +
%         \frac{1}{2}
%         \rbr{
%             0 + 0
%         } 
%         % =\frac{k}{4}
%         =\frac{k}{2}
%         .
%     \] 
%     Now consider when competitors are untruthful. 
%     In this setting, the optimal truthful policy will
%     not always 
%     give $\competitorb$ competitor
%     $\competitora$'s point. This is because doing so will always lead to a welfare of no more than 0,
%     thus such a policy is no better than one which exchanges zero resources.
%     Now consider a policy which only gives $\competitorb$ the point based on what is reported as $\valaii{1}$.
%     In this case $\competitora$ is incentivized to always report the value for $\valaii{1}$ for which $\competitorb$
%     does not receive the point as it increases their utility. However, this leads to a negative welfare.

%     Therefore, the optimal policy when competitors are untruthul recommends no exchange so 
%     \(
%         \EE\sbr{
%             \rewardss{
%                 \stopstateforpolicy{
%                     \opttruthexpolicy_k
%                 }
%             }
%         }
%         =0
%     \)
%     and thus
%     \[
%         \EE\sbr{
%             \rewardss{
%                 \optstrategytt{
%                     \optstoptime
%                 }
%             }
%         }
%         -
%         \EE\sbr{
%             \rewardss{
%                 \opttruthstrategytt{
%                     \opttruthstoptime
%                 }
%             }
%         }
%         \xrightarrow{k\rightarrow\infty}\infty .
%     \]
% \end{proof}  

\mycomment{before stating this result and when introducing truthfulness,
spend time explaining how the opt alg isn't truthful. If we put this in a section
4, then I can reference the example we walked through to claim optimality}

\mycomment{could add a figure of 2D stopping times which could
take up some space}
